%% GENERATED by scripts/paper/build_paper.py from docs/analysis/register/findings.yaml -- do not edit by hand;
%% rewrite the prose in main.tex or change the register / scripts/paper/paper_text.json and rebuild.

\subsection{You trust what you cannot check}\label{sec:qual:trust-unchecked}

\regtext{\textbf{Summary.} The most articulate operators describe trusting the Tactical Assistant \emph{because} they could not assess its plan, while the strategy card was simple enough to judge at a glance, and so was judged, often harshly. Under time pressure checking shrinks to a glance, and people say so with unusual candour: ``if we look at the recording I was doing exactly that''. A minority did check, and their trust was earned by comparing the assistant's plan with their own.} \par\noindent{\small\textit{Strength of evidence (register rating): qualitative. Your verdict: keep.}}

\subsubsection{Trusted because it cannot be checked}\label{sec:qual:trust-unchecked:cant-check}

\noindent \passtag{for the paper} \usernote{Your note: A nice few quotes that we can use to make a point.}

\noindent \regtext{The strategy card shows a few numbers per drone type beside the mission's own requirements, so its oddities are visible. The tactical plan is a web of routes on a map. P22 puts the consequence most clearly: they take the tactical plan ``at its word'' because they cannot see how far it is from optimal, and they suspect that is the only reason it seems so good. P19 says rating it would be ``disingenuous'' because they never really looked, and P01 that it is ``a lot harder to mentally challenge'' than drone numbers. This is the participants' own account of two quantitative results: the trust gap between the tiers (Section~\ref{sec:res:tier-trust}) and the visible versus invisible flaws (Section~\ref{sec:res:flaw-odd}). It suggests the gap is partly a gap in \emph{checkability}, not in quality.}

\paragraph*{What participants said.} \regtext{The clearest voices are P22, P19 and P01; P28 gives the same split from the other side (the strategic decision is the one you can analyse).}

\paragraph*{Quotes selected for this section.}

\begin{quotelist}

\item ``The tactical assistant, I never really saw it do anything wrong, but that's basically because I think the problem is so complex [that] seeing a more optimal solution is really hard. So I wonder if it's one of those things where these are based on the fact that I actually can't perceive the difference between its performance and the optimal.''\,---\,P22 (talk during the questionnaire survey1 0:14:45)

\item ``I will be honest with you. I don't think I was paying as much attention to the suggestions of the tactical assistant as I could have been. Just purely because the tasks were often very multi-drone. I think I would have committed a lot of extra thought-space to assigning all of the drones individually to those. So I definitely inadvertently trusted the tactical [agent]. I think it's worth saying any assessment of the tactical system would probably be disingenuous just because of my lack of engagement with it. But the strategic assistant, I think on the most part, I have absolutely no problems with. There was that one issue, but that's okay because I noticed it because the display is straightforward and it's clear enough that you can as a human spot discrepancies if they arise.''\,---\,P19 (talk during the questionnaire survey1 0:17:06)

\item ``i felt like i took a bigger pinch of salt for the strategic one because i felt like i could assess the quality of it more immediately whereas with the tactical assistant, I couldn't accurately and rapidly assess how well it was performing. So I just kind of said, OK, I have to take your word for it, essentially.''\,---\,P22 (talk during the questionnaire survey2 0:26:53)

\end{quotelist}

\clearpage

\subsubsection{Checking shrinks to a glance}\label{sec:qual:trust-unchecked:glance}

\noindent \passtag{extra: may make it in} \usernote{Your note: Less significant but might be worth a mention}

\noindent \regtext{Nobody describes careful checking under pressure. They describe a quick look (``does anything look ridiculous?''), knowing that a plan can look fine and not be, and accepting that trade because penalty accrues while they check. P31 gives the whole arc: over-checking at the start, under-checking by the end, less oversight as the session goes on, and a rating corrected against their own behaviour. P11 makes the expected-value argument explicitly: on a big mission even a dodgy plan pays, because of the time saved for the next mission. P28 is the exception, saying they checked every answer, though ``a quick browse''.}

\paragraph*{What participants said.} \regtext{P31 is the richest single source here; P29, P15 and P22 say it more briefly; P28 dissents.}

\paragraph*{Quotes selected for this section.}

\begin{quotelist}

\item {\upshape\footnotesize Researcher: Am I correct in saying you should be using the agents a little bit more often now?}\par ``Whatever lack of optimality that comes from using the AI agents is more than compensated for in their speed. Or at least in their speed of decision making. I also feel like the longer that this session has gone on, the less oversight that I've been giving it, and the more I've been focusing on the strategic layer and the incoming missions.''\,---\,P31 (think-aloud s2 0:20:58)

\item {\upshape\footnotesize Researcher: So you say you especially trust it with a large mission. You mean you're happy to rely on it because of the speed?}\par ``Yeah. I think even if the tactical agent gave me a dodgy suggestion and the mission failed and I lost points, I probably would have gained more points from the time saved''\,---\,P11 (debrief 0:29:11, includes the researcher's words)

\item ``Speed is the name of the game. So sometimes in these ones, I guess making a decision earlier is beneficial because at some point the other ones are going to pop up.''\,---\,P22 (think-aloud s1 0:04:20, includes the researcher's words)

\end{quotelist}

\clearpage

\subsubsection{Trust earned by comparison, and doubts that came with practice}\label{sec:qual:trust-unchecked:earned}

\noindent \passtag{for the paper} \usernote{Your note: Minimal but nice}

\noindent \regtext{Where trust was earned, it was by comparison. P32 trusted the tactical plan because it ``seemed to be doing what I was doing'' and the cards did not; P29 tested themself against the assistant early and found their own builds were the ones that failed; P11 decided they had been ``too harsh'' once Suggest matched their simple plan at the same speed. Practice cut both ways: P15 only saw the assistant's mistakes once they had time to look, and suspects they were there all along.}

\paragraph*{What participants said.} \regtext{P32, P29, P11 and P15.}

\paragraph*{Quotes selected for this section.}

\begin{quotelist}

\item {\upshape\footnotesize Researcher: So I would start with just how did you find that? Anything immediately that you would comment on?}\par ``I think I developed a stronger opinion on how much to trust each of them. By, I guess, comparing them to what I would have done. And I mean, I understood the task more as I went. So the tactical one seemed to be doing what I was doing. So quite reliably. And obviously the strategic one wasn't. So I guess it was very much comparing them to what I'd do. And then letting the tactical one get on with it.''\,---\,P32 (debrief 0:23:45)

\item {\upshape\footnotesize Researcher: The degree to which you trusted the decisions the AI was making, uh, and how that progressed throughout the study?}\par ``I mostly trusted it, although, uh, I did kind of want to test it more at the beginning, like comparing it, like doing a few manual ones myself to see, like, whether it would be better, faster. I guess it was kind of hard to keep track of which tasks were being finished when, but I noticed that the ones that failed were normally cause I'd just done a manual one. So I guess that made me want to use the AI more often.''\,---\,P29 (debrief 0:25:59)

\end{quotelist}

\clearpage

% q-process (Trust in the process, not the answer: what operators thought the assistants were) left out: your verdict this pass was prune.

\subsection{Who decides, and who answers for it}\label{sec:qual:who-decides}

\regtext{\textbf{Summary.} Asked about real search and rescue, nobody rejected assistants like these, but almost everyone wanted a human to keep the decision, and the reasons were moral as much as practical: accountability for a life not saved, a human ``connected'' to the operation, and the plan held in someone's head. People were candid that the human they want in the loop may not be able to check much. Against the grain, P13 wants enough automation that ``the end users could be anyone''.} \par\noindent{\small\textit{Strength of evidence (register rating): qualitative. Your verdict: keep.}}

\subsubsection{Accountability that just isn't there}\label{sec:qual:who-decides:accountability}

\noindent \passtag{for the paper}

\noindent \regtext{P31 sets out the dilemma: even if an AI saves more lives on average, when it fails to save one, ``that accountability just isn't there''. P22 weighs the same discomfort against saving five people sub-optimally rather than none, and settles on a human making the final decision. (P22's ``a bad approximation is better than no response at all'' refers back to an item of the pre-study questionnaire.) P31 also raised, unprompted, that this research ``will inevitably lead to it being used in military''.}

\paragraph*{What participants said.} \regtext{P31 and P22.}

\paragraph*{Quotes selected for this section.}

\begin{quotelist}

\item {\upshape\footnotesize Researcher: But one of my questions is if in five years' time this was much more sophisticated, far more complicated, how would you feel about search and rescue scenarios being conducted with AI systems of this type in this kind of setup?}\par ``Yeah, I feel like that's a very interesting question for more of a moral reason, because if you look at it entirely from the perspective of your option A, option B, and option A saves more lives, that always seems like a good thing, but taking it ... in a slightly less simplified way, you've got to make sort of moral judgments about how these resources are being allocated, and if someone's life was like failed to be saved, under a completely human system, you'd have some sort of like, not responsibility, what's the word I'm looking for? [interviewer suggests the word 'accountability']. Accountability, thank you! Of who was on the job that failed to save your life, whereas if it's an AI system, even if it performs on average better than a human, that accountability just isn't there.''\,---\,P31 (debrief 0:27:53, includes the researcher's words)

\item {\upshape\footnotesize Researcher: How do you feel about people using systems like this in, um, high risk scenarios?}\par ``I think definitely with emergency scenarios, it's best to just get something going, even if it's not the most perfect plan. And AI lets you do that quite good, like quite quickly.''\,---\,P29 (debrief 0:27:05, includes the researcher's words)

\end{quotelist}

\clearpage

\subsubsection{A human in the loop, even one who cannot check}\label{sec:qual:who-decides:loop}

\noindent \passtag{for the paper}

\noindent \regtext{Some want the human there regardless of what they can do: P24, pressed on whether a busy human could check, concedes they could not and still wants ``a human stand here''; P28 says ``humans should make that decision'' because ``we cannot guarantee what type of answers it will provide''. P32 gives the organisational version: measure where people and agents add value and assign the tiers accordingly.}

\paragraph*{What participants said.} \regtext{P24, P28 and P32.}

\paragraph*{Quotes selected for this section.}

\begin{quotelist}

\item {\upshape\footnotesize Researcher: If we took this five years in the future, and systems of this type were being used in real search and rescue scenarios with more complicated AIs and operators being assisted, how do you feel about that? And do you see any potential for good and for bad of doing that sort of thing?}\par ``I think if I were being organisationally strategic about that, it seems like it's really helpful. Assuming that you're human constrained (if you've got loads of spare humans, I'd want to see if the AI was actually better) but assuming that you've got less humans than you ideally would have, it feels like a good accelerator. I'd want the analytics on, I guess, like drone utilisation and mission success. I feel like that would be a good place to be making these decisions. And should I be, for example, in this case where I think the strategic one was less good than the tactical, if I were being really strategic, I'd be saying: 'okay, people are going to do the strategic stuff, and we're just going to trust it for the tactical'. So I guess very much doing that optimisation around: 'is it the agents or the humans who are adding the most value, and where do they add the most value', and then have them allocate them... I'd want to understand, well, just because I like to understand everything, I'd want to understand empirically actually how it was performing, who should be doing what.''\,---\,P32 (debrief 0:26:10, includes the researcher's words)

\end{quotelist}

\clearpage

\subsubsection{Ownership, and the opposite view}\label{sec:qual:who-decides:ownership}

\noindent \passtag{for the paper} \usernote{Your note: This is one of the better ones}

\noindent \regtext{For some the reason to keep work is not safety but connection. P11 does not want to be ``reduced to pressing allocate, suggest, deploy, because otherwise I might as well be in the next room'', and admits to pride and ``some sort of grudge'' against the computer. P21 keeps both decisions because the allocation already carries a plan in their head that the assistant would ``disrupt''. P13 takes the opposite view: automate everything not ethically constrained, so that the operator's skill matters less.}

\paragraph*{What participants said.} \regtext{P11 and P21 against P13.}

\paragraph*{Quotes selected for this section.}

\begin{quotelist}

\item ``I don't want to be reduced to pressing allocate, suggest, deploy, because otherwise I might as well be in the next room. And I think the small missions give me the chance to do that more. And I'm like, 'oh, I can just do that. Don't worry. I'll do that'. Done. And I feel satisfaction.''\,---\,P11 (debrief 0:30:55)

\item ``it would be a dreadful task without any of the tools, especially on the tactical side. It's just a lot of laborious dragging, holding shift at the right time.... And I guess I do feel I want to use (empahsis)my(/emphasis) solution, and I don't want to like, oh, you know, of course the computer thinks it's better. Or like personify the computer and feel some sort of grudge to hold it.''\,---\,P11 (debrief 0:26:55, includes the researcher's words)

\end{quotelist}

\clearpage

\subsection{Failure without feedback}\label{sec:qual:failure}

\regtext{\textbf{Summary.} Drone failures were random and uniform, and nothing on screen said so or said how likely they were. Operators felt the lack: P13 is ``making blind decisions'', P31 cannot tell whether their nudging helped, and several wanted failure rates tied to drone type or place built into the allocation. Without information they built their own theories (hostile places, ``fires or something'') and blamed themselves.} \par\noindent{\small\textit{Strength of evidence (register rating): qualitative. Your verdict: keep.}}

\subsubsection{Blind decisions, and failure data built into the plan}\label{sec:qual:failure:blind}

\noindent \passtag{for the paper}

\noindent \regtext{P13 treats allocation as a decision about risk that cannot be made well without the failure rate, and later proposes feeding a terrain-based estimate into the allocation rule so that the reserve is sized to the risk rather than padded. P10 wants failure rates by drone type from historic data. (In P11's and P10's debriefs the idea of per-area risk was raised by the researcher; P13 arrived at it themselves.)}

\paragraph*{What participants said.} \regtext{P13, P10 and P11.}

\paragraph*{Quotes selected for this section.}

\begin{quotelist}

\item ``I think in a in a real world scenario where if I imagine I'm a drone swarm operator you will have some numbers presumably about 'so what's the likelihood for drone failure' right? Based on historic data, and that will presumably bias (or inform, not bias) but inform your selection if I'm like okay the lifter drone is much heavier therefore there's more wear and tear on it or whatever it's more likely to fail down that means that my redundancy strategy might change from 'send one additional drone for the lifter' towards 'send 20 percent more more lifters than you need''\,---\,P10 (debrief 0:28:53, includes the researcher's words)

\end{quotelist}

\clearpage

\subsection{The hierarchy as one plan: how the two decisions connect}\label{sec:qual:hierarchy}

\regtext{\textbf{Summary.} Operators did not experience two independent assistants. They describe one plan that starts at the strategic step: over-allocate at the top and the shortage cascades down; allocate generously and there is ``really only one way'' to assign the drones; build the allocation yourself and the tactical plan is already in your head. Several place responsibility for failures with the strategic tier, and use the cross-mission view the cards lack (who can wait, what to hold back). Whether one plan held across the tiers is the right ideal is the question the practitioner's account, at the end, takes up and answers differently.} \par\noindent{\small\textit{Strength of evidence (register rating): qualitative. Your verdict: core.}}

\subsubsection{The strategic step carries the risk}\label{sec:qual:hierarchy:cascade}

\noindent \passtag{for the paper}

\noindent \regtext{P19 reaches the conclusion from the results: the tactical assistant ``did everything'', and the only remedy for failures is sending enough drones at the start, ``the strategic assistant's prerogative''. P28 weighs the same step from the other side: they hold drones back in the base because ``I cannot predict when the next task will appear'', and accept that a colour sent as a single drone may fail.}

\paragraph*{What participants said.} \regtext{P19 and P28.}

\paragraph*{Quotes selected for this section.}

\begin{quotelist}

\item ``Basically, it did everything. If failure rate of drones is unrelated to my assignment, so rate of failure as opposed to how many drones I send out, then all things considered, the tactical assistant is probably the best. The tactical assistant completed all the jobs and we didn't end up in any kind of a situation where we didn't have enough drones if it hadn't been for random drone failures. And the only way to remedy drone failures is to have sent enough drones at the beginning, which I think is the strategic assistant's prerogative.''\,---\,P19 (talk during the questionnaire survey1 0:18:53)

\item {\upshape\footnotesize Researcher: You found that one more demanding?}\par ``I'm trying to be more defensive in the strategy, so I'm trying to reserve more and more drones in the base, rather than send them out to finish the task quickly, because I cannot predict when the next task will appear. So luckily, there's only one task is holding for a long time, only one there, and another problem is [that] such a choice will make the drones in the case becomes not sufficient for that case. So if I only send out one drone for one specific color, and it fails, then I have to abort that mission, and I only abort one of them, because that happened, and the rest is still able to finish, so I think this, such a deployment is still all right.''\,---\,P28 (talk during the questionnaire survey2 0:28:05)

\end{quotelist}

\clearpage

\subsubsection{One plan across both decisions}\label{sec:qual:hierarchy:one-plan}

\noindent \passtag{for the paper}

\noindent \regtext{The card-then-Suggest pattern in the logs (Section~\ref{sec:res:size-decides}) is described from both sides. P10: a generous card means ``there is really only one way of assigning them, which means if I just hit the suggest button then I can focus on something else''. P21: building the allocation by hand lays out a plan in the mind, so switching to Suggest would disrupt it; their own plans chain drones through each other's tasks to get ``redundancy for free''. P21 also wonders whether the tactical assistant follows the strategic one's patterns. The practitioner's account at the end disagrees with this ideal: in staff planning each cell sees only its part, which is fine if each part has been through a sound process.}

\paragraph*{What participants said.} \regtext{P10 and P21.}

\paragraph*{Quotes selected for this section.}

\begin{quotelist}

\item ``it's better to use the autonomy on the tactical level on top of the strategic''\,---\,P10 (debrief 0:25:46, includes the researcher's words)

\item ``[Having manually set a strategic allocation] I prefer to [tactically] set the drones myself. Because I'm not sure how exactly it will be allocated [if I use the agent]. Because I have something in mind when I consider the allocation... When I make the selection for the drones, I kind of lay out a plan in my mind initially. So I feel like it's going to disrupt my thought process if I try to go around that.''\,---\,P21 (think-aloud s1 0:14:16, includes the researcher's words)

\item ``Actually, they might work better if I use the AI agent for the selection and that as well, because. Yeah, because when I assume when it makes a selection, when it tries to make the drone allocation, it follows certain patterns, which could then be adopted by the tactical agent afterwards. And if I try to think of something else, so I'm going to split them like that and change them in a certain way. I'm not sure if the tactical agent will be able to match my thought process.''\,---\,P21 (talk during the questionnaire between1 0:20:36)

\end{quotelist}

\clearpage

\subsubsection{The cross-mission view the cards lack}\label{sec:qual:hierarchy:cross-mission}

\noindent \regtext{The strategy cards look at one mission. Operators override them to act on what the cards cannot see: P22 strikes a ``middle ground'' to keep drones back for other tasks; P24 trims cards not because they are better than the AI but to save drones, and only then learns the cards see one mission; P19 reasons about who can afford to wait; P31 goes below even the Conservative card when other tasks are pending.}

\paragraph*{What participants said.} \regtext{P22, P24, P19 and P31.}

\paragraph*{Quotes selected for this section.}

\begin{quotelist}

\item {\upshape\footnotesize Researcher: Could you just summarize again, just in a sentence or two, how you were approaching each one, what scenarios you felt you were using each one?}\par ``I felt like sometimes I could find like a middle ground that had like maybe a little bit of gambling in there, essentially, that you wouldn't have the necessary coverage. But then also then that meant that you'd keep some more in reserve. So it wasn't just fully because it's cautious, it was also because it was actually to keep some in reserve for other tasks, because, yes, you're covering this task specifically very well. But that doesn't mean that then you actually have enough coverage for the, for the rest of your tasks. So I felt that often I could find, I could strike a middle ground. So I felt comfortable assessing and then overriding the strategic assistant.''\,---\,P22 (debrief 0:28:28)

\item ``And the problem is we don't know whether a extremely critical job is going to show up later.''\,---\,P19 (think-aloud s1 0:09:07)

\item {\upshape\footnotesize Researcher: Now that I've had a little bit more time, I'm just thinking if my opinions have changed on this slightly. And for what reason?}\par ``If you wanted to send as few drones as possible, you could send four fast drones, one lifter, one camera. Whereas even the conservative recommendation recommended sending two of the lifters, and the cameras in case of drone failure. When I was in that situation, I had a lot of other pending tasks, so I went more conservative than the conservative recommendation.''\,---\,P31 (talk during the questionnaire survey2 0:23:49, includes the researcher's words)

\end{quotelist}

\clearpage

% q-real-ops (What a real operation would need) left out: your verdict this pass was prune.

\subsection{New jobs for an assistant}\label{sec:qual:new-jobs}

\regtext{\textbf{Summary.} Asked what the assistants should do, people rarely asked for more automation of the same kind. They asked for assistants that check: one that reviews the operator's own plan for slips rather than taking it over (P11). They asked for smarter rules (redundancy scaled to failure risk, backup routes in the tactical plan), for options between and beyond the two cards, and for an interface that knows what they have already planned. Two (P27, P32) asked for something neither assistant does: help deciding which waiting mission to open next.} \par\noindent{\small\textit{Strength of evidence (register rating): qualitative. Your verdict: keep.}}

\subsubsection{An assistant for what to open next}\label{sec:qual:new-jobs:queue-order}

\noindent \passtag{for the paper}

\noindent \regtext{Neither assistant helps with the first decision of every allocation: which waiting mission to open. P32 asked for it in play (``a suggestion of what to do next would be handy''), because a mission the reserve cannot cover sits at the top of the list anyway; P27 gave a rule, missions that can be done first, then by criticality, allowing for what is already running. The logs show the gap such an assistant would fill. With two or more missions waiting (411 decisions, 33 people), operators open the oldest (49\%; chance 34\%) and lean towards the more critical ($\times$1.45 per category above the waiting set's average, p = .02), but not towards missions they can actually cover: whether the reserve could supply a mission's drones makes no difference ($\times$1.04, p = .79), and at 41\% of decisions the oldest waiting mission could not be covered. P27's order differs from what operators did at 59\% of decisions. The queue had a ``Sort: Score'' toggle; the researcher told P27 only one participant had found it, and it is not logged.}

\paragraph*{Key numbers.}

\begin{itemize}[leftmargin=1.2em,itemsep=1pt]

\item Decisions with 2+ missions waiting: 411 (33 people). Opened the oldest 49\% (chance 34\%); when one waiting mission was uniquely the most critical, opened it 49\% (chance 33\%)

\item Choice among the waiting missions (conditional logit, person bootstrap, 400 draws): oldest $\times$2.75 [1.97, 4.12], p < .001; newest $\times$1.64 [1.17, 2.32], p < .001; criticality, per category above the set's mean $\times$1.45 [1.09, 2.15], p = .02; coverable from the reserve $\times$1.04 [0.72, 1.58], p = .79; size per 3 drones $\times$0.95 [0.78, 1.11], p = .61

\item The oldest waiting mission could not be covered from the reserve at 41\% of these decisions; some waiting mission could not be covered at 57\%. When some could not, the one opened was coverable 32\% of the time (chance 39\%)

\item P27's order (coverable, then most critical, then oldest) picks a different mission from oldest-first at 59\% of decisions; there operators opened P27's pick 35\% of the time (chance 30\%)

\end{itemize}

\paragraph*{Method.} {\small Same decisions and waiting sets as the triage choice model (load.py), with two features added per waiting mission: its category (A--E, the penalty rate the operator sees), centred on the set; and whether the drones available in the reserve just before the decision covered its primary composition in every colour. Coverable is strict: tasks also have substitute compositions, and drones returning to the hub are not counted, so it understates what an operator could have done. scripts/register/pass2/queue\_order.py; results in pass2/load.json, queueOrder.}

\paragraph*{What participants said.} \regtext{P27 and P32; both cuts include the researcher's words (flagged).}

\clearpage
